
\section{基金与项目申请经历}
% increase linespacing [parsep=0.5ex]
\datedsubsection{\textbf{COVID-19在通勤过程中的传播建模}, \$90,000}{2020 -- 2021}
\begin{itemize}[parsep=0.5ex]
  \item \textbf{导师为PI，我为第一完成人，负责撰写基金申请和设计研究框架}
  \item 项目旨在为MIT社区开发COVID-19传播风险监控系统，我们主要研究人们在通勤过程中的感染风险建模。项目获得MIT Quest for Intelligence约9万美金资金支持，为期一年。
\end{itemize}
% \datedsubsection{\textbf{个性化路径推荐系统设计}, \$150,000 }{未中}
% \begin{itemize}[parsep=0.5ex]
%   \item \textbf{导师为PI，我为第一完成人，负责撰写基金申请和设计研究框架}
%   \item 与Ferrovial公司合作，Ferrovial公司在美国范围内拥有大量的高速公路收费设施，我们希望为他们的用户设计个性化路径推荐算法和奖励机制，可以降低人们的出行时间，提供公路使用效率，同时提高公司收入，目前项目正在申请中。
% \end{itemize}
\datedsubsection{\textbf{基于区块链的绿色出行奖励生态设计}, \$200,000}{未中}
\begin{itemize}[parsep=0.5ex]
  \item \textbf{导师为PI，我和两位合作者为共同第一完成人，负责撰写基金申请和设计研究框架}
  \item 该项目正在申请美国交通部小企业创新研究计划基金（US DOT Small Business Innovation Research Program）。我们希望设计基于区块链的出行奖励机制，通过发放虚拟货币，鼓励人们进行绿色出行。我们设计的货币分发算法可以保证虚拟货币的最低价值（与减少的碳排放挂钩），未来目的是建立基于区块链的碳权交易市场。
\end{itemize}
% \datedsubsection{\textbf{公共交通与COVID-19的传播：模型、控制和管理}, \$200,000}{未中}
% \begin{itemize}[parsep=0.5ex]
%   \item \textbf{导师为PI，我为第一完成人，负责撰写基金申请和设计研究框架}
%   \item 该项目旨在理解COVID-19在公共交通系统（公交车、地铁）中传播的机制，提出相应的公共交通系统运行与制策略，以应对当前的疫情，并为未来其他传染疾病做准备。我负责撰写基金申请，设计研究框架，该项目曾申请Greater Boston Consortium on Pathogen Readiness Evergrande COVID-19 Response Fund Awards，但是未获接收。
% \end{itemize}
%% Reference
%\newpage
%\bibliographystyle{IEEETran}
%\bibliography{mycite}